\documentclass[11pt]{article}
\usepackage[margin=1in]{geometry}
\usepackage{amsmath,amssymb,amsthm}
\usepackage{hyperref}

\newtheorem{theorem}{Theorem}
\newtheorem{corollary}{Corollary}
\newtheorem{remark}{Remark}

\newcommand{\K}{\mathbb K}
\newcommand{\T}{\mathbb T}
\newcommand{\conv}{\operatorname{conv}}
\newcommand{\Spear}{\operatorname{Spear}}
\newcommand{\ext}{\operatorname{ext}}
\newcommand{\Rea}{\operatorname{Re}}

\title{Finite-rank aDP operators are lush}
\author{Run fa\_banach\_001, lane 19}
\date{2026-06-16}

\begin{document}
\maketitle

\section*{Source and target}

The source paper is Kadets, Mart\'in, Mer\'i, and P\'erez,
\emph{Spear operators between Banach spaces}, arXiv:1701.02977.
All Banach spaces below are real or complex.

After Proposition 6.5, the source paper says that it is not known whether
the finite-dimensional codomain theorem, or part of it, remains true when
only the range of the operator is finite-dimensional.  Chapter 10 also asks
whether the rank-one results extend to finite-rank operators.  The result
below proves the aDP/spear/lush equivalence part for finite-dimensional
range, hence for finite-rank operators.

\section*{Definitions and source facts}

A norm-one operator \(G:X\to Y\) has the alternative Daugavet property
(aDP) if the alternative Daugavet equation holds for every rank-one
operator \(T:X\to Y\).  A spear operator satisfies the analogous equation
for every bounded operator \(T:X\to Y\), and every lush operator is spear.

We use two facts from the source paper.

\begin{enumerate}
\item Theorem 3.2(iv): if \(G:X\to Y\) has the aDP, then for every
\(x^*\in X^*\) the set
\[
D(x^*)=\left\{y^*\in \ext B_{Y^*}:
\|G^*y^*+\T x^*\|=1+\|x^*\|\right\}
\]
is dense \(G_\delta\) in \((\ext B_{Y^*},w^*)\).
\item Proposition 3.32(a): if
\[
\left\{y^*\in B_{Y^*}:G^*y^*\in \Spear(X^*)\right\}
\]
is norming for \(Y\), then \(G\) is lush.
\end{enumerate}

\section*{Proof Intuition}

For one test vector \(x^*\in X^*\), the aDP gives many extreme functionals
\(y^*\) for which \(G^*y^*\) satisfies the spear equation against that
particular \(x^*\).  For finitely many test vectors, Baire category gives
one \(y^*\) satisfying all finitely many equations inside any desired
weak-star slice.

The finite-rank assumption lets us pass from finitely many equations to all
equations.  Indeed, the set \(G^*(B_{Y^*})\) sits in a finite-dimensional
subspace of \(X^*\), so weak-star compact slices of \(B_{Y^*}\) have
norm-compact images under \(G^*\).  The closed sets expressing each spear
equation have the finite-intersection property, hence a common point.  This
common point is a spear functional in \(X^*\), and because the weak-star
slice was arbitrary, such functionals norm \(Y\).  Proposition 3.32(a) then
turns this norming family into lushness of \(G\).

\section*{Main theorem}

\begin{theorem}
Let \(X,Y\) be Banach spaces and let \(G:X\to Y\) be a norm-one
finite-rank operator.  If \(G\) has the alternative Daugavet property, then
\(G\) is lush.
\end{theorem}

\begin{proof}
It suffices, by Proposition 3.32(a) of the source paper, to prove that
\[
\mathcal N=\{y^*\in B_{Y^*}:G^*y^*\in \Spear(X^*)\}
\]
is norming for \(Y\).

Fix \(y_0\in S_Y\) and \(0<\varepsilon<1\).  Let
\[
U=\{y^*\in B_{Y^*}:\Rea y^*(y_0)>1-\varepsilon\}.
\]
The set \(U\cap \ext B_{Y^*}\) is nonempty.  Indeed, \(U\) is a nonempty
weak-star slice of \(B_{Y^*}\), and if every extreme point of \(B_{Y^*}\)
lay outside \(U\), then the weak-star closed convex hull of the extreme
points would also lie outside \(U\), contradicting Krein--Milman.

For each \(x^*\in X^*\), Theorem 3.2(iv) gives a dense \(G_\delta\) subset
\[
D(x^*)=\{y^*\in \ext B_{Y^*}:
\|G^*y^*+\T x^*\|=1+\|x^*\|\}
\]
of \((\ext B_{Y^*},w^*)\).  Hence, if \(F\subset X^*\) is finite, then
\[
U\cap \ext B_{Y^*}\cap \bigcap_{x^*\in F}D(x^*)\ne\varnothing .
\]

Now put \(E=G^*(Y^*)\).  Since \(G\) has finite rank, \(E\) is a
finite-dimensional subspace of \(X^*\).  Let
\[
K=G^*(\overline U^{\,w^*})\subset E.
\]
This set is compact in norm, because \(\overline U^{\,w^*}\) is weak-star
compact and \(E\) is finite-dimensional.  For each \(x^*\in X^*\), define
\[
K_{x^*}=\{e\in K:\|e+\T x^*\|=1+\|x^*\|\}.
\]
Each \(K_{x^*}\) is closed in \(K\).  The preceding Baire-category
observation says that the family \(\{K_{x^*}:x^*\in X^*\}\) has the
finite-intersection property: for finitely many \(x_1^*,\ldots,x_n^*\), pick
\(y^*\in U\cap \ext B_{Y^*}\cap\bigcap_{j=1}^nD(x_j^*)\), and then
\(G^*y^*\in\bigcap_{j=1}^nK_{x_j^*}\).

By compactness of \(K\), there exists \(e_0\in\bigcap_{x^*\in X^*}K_{x^*}\).
Since \(e_0\in K=G^*(\overline U^{\,w^*})\), choose
\(y_0^*\in\overline U^{\,w^*}\) with \(G^*y_0^*=e_0\).  The equality for
\(x^*=0\) gives \(\|e_0\|=1\), and the equality for every \(x^*\in X^*\)
gives
\[
\|e_0+\T x^*\|=1+\|x^*\|\qquad (x^*\in X^*).
\]
Thus \(e_0=G^*y_0^*\in \Spear(X^*)\).  Also
\(\Rea y_0^*(y_0)\ge 1-\varepsilon\), because
\(y_0^*\in \overline U^{\,w^*}\).  Therefore \(\mathcal N\) is norming for
\(Y\).  Proposition 3.32(a) now implies that \(G\) is lush.
\end{proof}

\begin{corollary}
Let \(G:X\to Y\) be a norm-one finite-rank operator.  Then the following
are equivalent:
\[
G\text{ is lush},\qquad
G\text{ is a spear operator},\qquad
G\text{ has the aDP}.
\]
\end{corollary}

\begin{proof}
Every lush operator is spear, and every spear operator has the aDP.  The
theorem proves the converse implication from aDP to lush for finite-rank
operators.
\end{proof}

\section*{Consequences and limitations}

This theorem strengthens the earlier lane-19 norm-attainment packet:
finite-rank aDP operators attain their norm because they are lush, and the
source paper proves that lush operators attain their norm.

The theorem does not settle the finite-rank analogue of Proposition 7.8 of
the source paper.  In other words, it remains open from this argument
whether \(G^*:Y^*\to X^*\) must be lush whenever \(G:X\to Y\) is finite-rank
and has the aDP.

\section*{Novelty notes}

The proof uses only source-paper tools plus finite-dimensional compactness
of \(G^*(B_{Y^*})\).  Cheap run-index searches found no prior packet for this
finite-rank aDP/lushness equivalence.  Local checks against arXiv:2306.02645
and arXiv:2306.02638 did not find the statement.  Web searches on
2026-06-16 for exact combinations of ``finite-rank'', ``alternative Daugavet
property'', ``finite-dimensional range'', ``spear operators'', and ``lush''
found no exact duplicate.

\begin{thebibliography}{9}

\bibitem{KMMP}
V. Kadets, M. Mart\'in, J. Mer\'i, and D. P\'erez,
\emph{Spear operators between Banach spaces},
arXiv:1701.02977.

\end{thebibliography}

\end{document}

